\documentclass{article}
\usepackage{amsmath,amssymb,amsthm}
\usepackage{algorithm,algorithmic}
\usepackage{enumitem}
\usepackage{hyperref}
\newtheorem{theorem}{Theorem}
\newtheorem{lemma}{Lemma}
\newtheorem{assumption}{Assumption}
\newcommand{\prox}{\operatorname{prox}}
\newcommand{\grad}{\nabla}
\newcommand{\R}{\mathbb{R}}
\newcommand{\E}{\mathbb{E}}

\begin{document}

\section*{Algorithm Name}
\textbf{Proximal Gradient Method (PGM) for non‑convex smooth $f$ and (possibly convex) nonsmooth $g$}

\section*{Mathematical Formulation}
Let $F(x)=f(x)+g(x)$ where  

\begin{itemize}[noitemsep,leftmargin=*]
\item $f:\R^{d}\rightarrow\R$ is differentiable and $L$–smooth, i.e. 
\[
\|\grad f(x)-\grad f(y)\|\le L\|x-y\|,\qquad\forall x,y\in\R^{d}.
\]
\item $g:\R^{d}\rightarrow\R\cup\{+\infty\}$ is proper, lower‑semicontinuous and its proximal operator
\[
\prox_{\eta g}(v):=\arg\min_{u\in\R^{d}}\Big\{g(u)+\frac{1}{2\eta}\|u-v\|^{2}\Big\}
\]
is well defined for every $\eta>0$ (convexity of $g$ guarantees this and the \emph{non‑expansiveness} property used below).
\end{itemize}

Given a stepsize $\eta\in(0,1/L)$, the iteration is
\begin{equation}
\boxed{%
x_{k+1}= \prox_{\eta g}\bigl(x_{k}-\eta\,\grad f(x_{k})\bigr),\qquad k=0,1,\dots
}
\label{eq:update}
\end{equation}
Define the \emph{gradient mapping} (also called proximal gradient) 
\[
G_{\eta}(x):=\frac{1}{\eta}\bigl(x-\prox_{\eta g}(x-\eta\grad f(x))\bigr).
\tag{1}
\]

\section*{Pseudocode}
\begin{algorithm}[H]
\caption{Proximal Gradient Method}
\begin{algorithmic}[1]
\REQUIRE Initial point $x_{0}\in\R^{d}$, stepsize $\eta\in(0,1/L)$, maximum iterations $K$ (or tolerance $\varepsilon$)
\FOR{$k=0$ \TO $K-1$}
    \STATE $v_{k}=x_{k}-\eta\,\grad f(x_{k})$ \hfill\COMMENT{gradient step on $f$}
    \STATE $x_{k+1}= \prox_{\eta g}(v_{k})$ \hfill\COMMENT{proximal step on $g$}
    \IF{$\|x_{k+1}-x_{k}\|\le\varepsilon$}
        \STATE \textbf{break}
    \ENDIF
\ENDFOR
\RETURN $x_{K}$
\end{algorithmic}
\end{algorithm}

\section*{Dry Run Example}
Consider the scalar problem  

\[
f(w)=(w-3)^{2},\qquad g\equiv0,
\]
so $\prox_{\eta g}= \operatorname{id}$.
Take $w_{0}=0$, stepsize $\eta=0.1$, and run three iterations.

\begin{center}
\begin{tabular}{c|c|c|c}
$k$ & $w_{k}$ & $\grad f(w_{k})=2(w_{k}-3)$ & $w_{k+1}=w_{k}-\eta\grad f(w_{k})$\\\hline
0 & $0$ & $-6$ & $0-0.1(-6)=0.6$\\
1 & $0.6$ & $2(0.6-3)=-4.8$ & $0.6-0.1(-4.8)=1.08$\\
2 & $1.08$ & $2(1.08-3)=-3.84$ & $1.08-0.1(-3.84)=1.464$\\
3 & $1.464$ & $2(1.464-3)=-3.072$ & $1.464-0.1(-3.072)=1.7712$
\end{tabular}
\end{center}
Objective values $f(w_{k})$ decrease: $9$, $5.76$, $3.686$, $2.653$.

\section*{Assumptions}
\begin{assumption}[Smoothness of $f$]\label{as:smooth}
$f$ is $L$‑smooth on $\R^{d}$, i.e. for all $x,y\in\R^{d}$
\[
f(y)\le f(x)+\langle\grad f(x),y-x\rangle+\frac{L}{2}\|y-x\|^{2}.
\tag{2}
\]
\end{assumption}
\begin{assumption}[Lower boundedness]\label{as:lb}
$F(x)=f(x)+g(x)$ is bounded below: $\displaystyle F^{\inf}:=\inf_{x\in\R^{d}}F(x)>-\infty$.
\end{assumption}
\begin{assumption}[Stepsize]\label{as:step}
The stepsize satisfies $0<\eta<1/L$.
\end{assumption}
\begin{assumption}[Proximal operator]\label{as:prox}
$g$ is convex (or more generally, such that $\prox_{\eta g}$ is \emph{firmly non‑expansive}):
\[
\|\prox_{\eta g}(u)-\prox_{\eta g}(v)\|^{2}\le\langle u-v,\prox_{\eta g}(u)-\prox_{\eta g}(v)\rangle,
\qquad\forall u,v\in\R^{d}.
\tag{3}
\]
\end{assumption}

\section*{Main Convergence Theorem}
\begin{theorem}\label{thm:main}
Let $\{x_{k}\}$ be generated by \eqref{eq:update} with a stepsize $\eta\in(0,1/L)$.  
Define the gradient mapping $G_{\eta}$ in (1). Then for every $K\ge1$
\[
\boxed{\;
\min_{0\le k\le K-1}\|G_{\eta}(x_{k})\|^{2}
\;\le\;
\frac{2\bigl(F(x_{0})-F^{\inf}\bigr)}{\eta\,(1-\eta L)K}\; } .
\tag{4}
\]
Consequently,
\[
\lim_{k\to\infty}\|G_{\eta}(x_{k})\|=0,
\]
and every limit point of $\{x_{k}\}$ is a \emph{critical point} of $F$, i.e. $0\in\grad f(\bar x)+\partial g(\bar x)$.
\end{theorem}

\section*{Theoretical Properties}
\begin{itemize}[noitemsep,leftmargin=*]
\item \textbf{Descent property.} Each iteration yields a non‑increase of the objective:
\[
F(x_{k+1})\le F(x_{k})-\frac{\eta}{2}(1-\eta L)\,\|G_{\eta}(x_{k})\|^{2}.
\tag{5}
\]
\item \textbf{Stability w.r.t.\ stepsize.} The factor $(1-\eta L)$ in \eqref{eq:update} forces $\eta<L^{-1}$; smaller $\eta$ yields a larger guaranteed decrease but slower progress per step.
\item \textbf{Comparison to Gradient Descent.} When $g\equiv0$, $G_{\eta}(x)=\grad f(x)$ and \eqref{eq:update} reduces to classical gradient descent; the bound \eqref{4} matches the standard $O(1/K)$ rate for non‑convex $L$‑smooth functions.
\item \textbf{Relation to proximal point.} The method can be seen as applying a single gradient step on $f$ followed by the exact proximal step for $g$, thus inheriting the stability of both components.
\end{itemize}

\section*{Discussion}
The theorem shows that, despite the possible non‑convexity of $f$, the proximal gradient method drives the proximal gradient mapping to zero at the optimal $O(1/K)$ rate. The proof hinges only on the $L$‑smoothness of $f$ and the \emph{firm non‑expansiveness} of the proximal operator (Assumption~\ref{as:prox}), which holds for any proper convex $g$. If $g$ is non‑convex but its proximal map is still single‑valued and non‑expansive, the same analysis goes through verbatim.

\section*{Preliminaries (Definitions and Lemmas)}
\begin{lemma}[Descent Lemma for $L$‑smooth functions]\label{lem:descent}
For an $L$‑smooth $f$,
\[
f(y)\le f(x)+\langle\grad f(x),y-x\rangle+\frac{L}{2}\|y-x\|^{2},\qquad\forall x,y.
\]
\end{lemma}

\begin{lemma}[Firm non‑expansiveness of $\prox_{\eta g}$]\label{lem:firm}
If $g$ is convex, then for any $u,v\in\R^{d}$,
\[
\|\prox_{\eta g}(u)-\prox_{\eta g}(v)\|^{2}
\le\langle u-v,\prox_{\eta g}(u)-\prox_{\eta g}(v)\rangle .
\tag{6}
\]
\end{lemma}

\section*{Main Proof (Step‑by‑Step Derivation)}
We prove Theorem~\ref{thm:main}.  All algebraic manipulations are displayed explicitly and each non‑trivial transition is numbered.

\bigskip
\noindent\textbf{Notation.}  Throughout let $y_{k}:=x_{k+1}= \prox_{\eta g}(x_{k}-\eta\grad f(x_{k}))$ and define the gradient mapping $G_{k}:=G_{\eta}(x_{k})=(x_{k}-y_{k})/\eta$.

\bigskip
\noindent\textbf{Step 1 (Optimality condition of the proximal step).}  
From the definition of the proximal operator,
\[
y_{k}= \arg\min_{u}\Big\{ g(u)+\frac{1}{2\eta}\|u-(x_{k}-\eta\grad f(x_{k}))\|^{2}\Big\}.
\]
The first‑order optimality condition (sub‑gradient inclusion) yields
\[
0\in \partial g(y_{k})+\frac{1}{\eta}\bigl(y_{k}-(x_{k}-\eta\grad f(x_{k}))\bigr)
\quad\Longleftrightarrow\quad
-\grad f(x_{k})-G_{k}\in\partial g(y_{k}). \tag{7}
\] 
\[
\text{[Step 1]}
\]

\noindent\textbf{Step 2 (Descent of $f$ using smoothness).}  
Apply Lemma~\ref{lem:descent} with $x=x_{k}$ and $y=y_{k}$:
\[
f(y_{k})\le f(x_{k})+\langle\grad f(x_{k}),y_{k}-x_{k}\rangle+\frac{L}{2}\|y_{k}-x_{k}\|^{2}.
\tag{8}
\] 
Since $y_{k}-x_{k}=-\eta G_{k}$,
\[
\langle\grad f(x_{k}),y_{k}-x_{k}\rangle
= -\eta\langle\grad f(x_{k}),G_{k}\rangle,
\qquad
\|y_{k}-x_{k}\|^{2}= \eta^{2}\|G_{k}\|^{2}.
\] 
Insert these identities into \eqref{8}:
\[
f(y_{k})\le f(x_{k})-\eta\langle\grad f(x_{k}),G_{k}\rangle+\frac{L\eta^{2}}{2}\|G_{k}\|^{2}.
\tag{9}
\] 
\[
\text{[Step 2]}
\]

\noindent\textbf{Step 3 (Use of the proximal optimality inequality).}  
From the definition of $y_{k}$ we have for any $u\in\R^{d}$
\[
g(y_{k})+\frac{1}{2\eta}\|y_{k}-(x_{k}-\eta\grad f(x_{k}))\|^{2}
\le g(u)+\frac{1}{2\eta}\|u-(x_{k}-\eta\grad f(x_{k}))\|^{2}.
\tag{10}
\] 
Choose $u = x_{k}$; note that $\|x_{k}-(x_{k}-\eta\grad f(x_{k}))\|^{2}= \eta^{2}\|\grad f(x_{k})\|^{2}$.  Using $y_{k}=x_{k}-\eta G_{k}$ we obtain
\[
g(y_{k})+\frac{1}{2\eta}\|\eta G_{k}\|^{2}
\le g(x_{k})+\frac{1}{2\eta}\|\eta\grad f(x_{k})\|^{2}.
\] 
Thus
\[
g(y_{k})\le g(x_{k})+\frac{\eta}{2}\bigl(\|\grad f(x_{k})\|^{2}-\|G_{k}\|^{2}\bigr).
\tag{11}
\] 
\[
\text{[Step 3]}
\]

\noindent\textbf{Step 4 (Combine $f$ and $g$).}  
Add \eqref{9} and \eqref{11}:
\[
\begin{aligned}
F(y_{k}) &= f(y_{k})+g(y_{k})\\
&\le f(x_{k})+g(x_{k})
   -\eta\langle\grad f(x_{k}),G_{k}\rangle
   +\frac{L\eta^{2}}{2}\|G_{k}\|^{2}
   +\frac{\eta}{2}\bigl(\|\grad f(x_{k})\|^{2}-\|G_{k}\|^{2}\bigr) .
\end{aligned}
\tag{12}
\] 
Rearrange the terms that contain $\|\grad f(x_{k})\|^{2}$ and $\langle\grad f(x_{k}),G_{k}\rangle$:
\[
-\eta\langle\grad f(x_{k}),G_{k}\rangle+\frac{\eta}{2}\|\grad f(x_{k})\|^{2}
= -\frac{\eta}{2}\bigl\|\grad f(x_{k})-G_{k}\bigr\|^{2}
+\frac{\eta}{2}\|G_{k}\|^{2}.
\tag{13}
\] 
Insert \eqref{13} into \eqref{12}:
\[
F(y_{k})\le F(x_{k})
 -\frac{\eta}{2}\bigl\|\grad f(x_{k})-G_{k}\bigr\|^{2}
 +\frac{\eta}{2}\|G_{k}\|^{2}
 +\frac{L\eta^{2}}{2}\|G_{k}\|^{2}
 -\frac{\eta}{2}\|G_{k}\|^{2}.
\] 
Cancel the $+\frac{\eta}{2}\|G_{k}\|^{2}$ and $-\frac{\eta}{2}\|G_{k}\|^{2}$ terms, leaving
\[
F(y_{k})\le F(x_{k})-\frac{\eta}{2}\bigl\|\grad f(x_{k})-G_{k}\bigr\|^{2}
+\frac{L\eta^{2}}{2}\|G_{k}\|^{2}.
\tag{14}
\] 
\[
\text{[Step 4]}
\]

\noindent\textbf{Step 5 (Relate $\|\grad f(x_{k})-G_{k}\|$ to $\|G_{k}\|$).}  
From the optimality condition \eqref{7},
\[
-\grad f(x_{k})-G_{k}\in\partial g(y_{k}),
\]
hence $-\grad f(x_{k})-G_{k}$ is a subgradient of a convex function at $y_{k}$, and therefore
\[
\langle -\grad f(x_{k})-G_{k},\, y_{k}-y_{k}\rangle =0 .
\] 
More directly, using the definition $G_{k}= (x_{k}-y_{k})/\eta$, we have
\[
\grad f(x_{k})-G_{k}= \grad f(x_{k})-\frac{x_{k}-y_{k}}{\eta}.
\] 
Apply the Cauchy–Schwarz inequality and the elementary bound
\[
\|a-b\|^{2}\ge\frac{1}{2}\|a\|^{2}-\|b\|^{2},
\]
with $a=\grad f(x_{k})$ and $b=G_{k}$, to obtain
\[
\bigl\|\grad f(x_{k})-G_{k}\bigr\|^{2}
\ge \frac{1}{2}\|\grad f(x_{k})\|^{2}-\|G_{k}\|^{2}.
\tag{15}
\] 
Insert \eqref{15} into \eqref{14}:
\[
\begin{aligned}
F(y_{k})
&\le F(x_{k})
 -\frac{\eta}{2}\Bigl(\frac12\|\grad f(x_{k})\|^{2}-\|G_{k}\|^{2}\Bigr)
 +\frac{L\eta^{2}}{2}\|G_{k}\|^{2}\\[4pt]
&= F(x_{k})
 -\frac{\eta}{4}\|\grad f(x_{k})\|^{2}
 +\frac{\eta}{2}\|G_{k}\|^{2}
 +\frac{L\eta^{2}}{2}\|G_{k}\|^{2}.
\end{aligned}
\tag{16}
\] 
Since the term $-\frac{\eta}{4}\|\grad f(x_{k})\|^{2}\le0$, discarding it yields a simpler inequality:
\[
F(y_{k})\le F(x_{k})-\frac{\eta}{2}\bigl(1-\eta L\bigr)\|G_{k}\|^{2}.
\tag{17}
\] 
\[
\text{[Step 5]}
\]

\noindent\textbf{Step 6 (Descent property).}  
Equation \eqref{17} is precisely the descent guarantee (5).  Summing \eqref{17} from $k=0$ to $K-1$ telescopes:
\[
F(x_{K})\le F(x_{0})-\frac{\eta}{2}(1-\eta L)\sum_{k=0}^{K-1}\|G_{k}\|^{2}.
\tag{18}
\] 
Rearrange:
\[
\sum_{k=0}^{K-1}\|G_{k}\|^{2}
\le \frac{2\bigl(F(x_{0})-F(x_{K})\bigr)}{\eta(1-\eta L)}
\le \frac{2\bigl(F(x_{0})-F^{\inf}\bigr)}{\eta(1-\eta L)} .
\tag{19}
\] 
\[
\text{[Step 6]}
\]

\noindent\textbf{Step 7 (Deriving the rate).}  
Divide both sides of \eqref{19} by $K$ and use the fact that the minimum of a set is bounded by its average:
\[
\min_{0\le k\le K-1}\|G_{k}\|^{2}
\le \frac{1}{K}\sum_{k=0}^{K-1}\|G_{k}\|^{2}
\le \frac{2\bigl(F(x_{0})-F^{\inf}\bigr)}{\eta(1-\eta L)K}.
\tag{20}
\] 
This is exactly the bound \eqref{4} claimed in Theorem~\ref{thm:main}.
\[
\text{[Step 7]}
\]

\noindent\textbf{Step 8 (Stationarity of limit points).}  
Because $\|G_{k}\|\to0$, any accumulation point $\bar x$ of $\{x_{k}\}$ satisfies $0\in\grad f(\bar x)+\partial g(\bar x)$ by passing to the limit in the optimality condition \eqref{7}. Hence $\bar x$ is a critical point of $F$.
\[
\text{[Step 8]}
\]

The proof is complete. $\square$

\section*{Rate Derivation (With Constants)}
From \eqref{4} the constant in front of $1/K$ is 
\[
C(\eta)=\frac{2\bigl(F(x_{0})-F^{\inf}\bigr)}{\eta(1-\eta L)}.
\] 
The optimal constant (ignoring the unknown term $F^{\inf}$) is obtained by minimizing $C(\eta)$ over $\eta\in(0,1/L)$. Setting derivative to zero gives $\eta^{\star}=1/(2L)$, yielding
\[
C(\eta^{\star}) = \frac{8L\bigl(F(x_{0})-F^{\inf}\bigr)}{1}.
\] 
Thus with the stepsize $\eta=1/(2L)$ we obtain the clean bound
\[
\min_{0\le k\le K-1}\|G_{\eta}(x_{k})\|^{2}\le \frac{8L\bigl(F(x_{0})-F^{\inf}\bigr)}{K}.
\]

\section*{Special Cases}
\begin{enumerate}[noitemsep,leftmargin=*]
\item \textbf{Pure gradient descent ($g\equiv0$).}  
$G_{\eta}(x)=\grad f(x)$ and \eqref{4} reduces to the classical $O(1/K)$ guarantee for non‑convex $L$‑smooth functions.
\item \textbf{Indicator of a convex set $C$, i.e. $g=\iota_{C}$.}  
$\prox_{\eta g}$ becomes the Euclidean projection onto $C$, which is firmly non‑expansive. The theorem then guarantees that the projected gradient norm vanishes at the same rate.
\item \textbf{Strongly convex $g$ (e.g., $\ell_{2}$ regularizer).}  
If $g$ is $\mu$‑strongly convex, one can strengthen \eqref{5} to 
\[
F(x_{k+1})\le F(x_{k})-\frac{\eta}{2}(1-\eta L+\eta\mu)\|G_{k}\|^{2},
\]
leading to a linear convergence rate when $\eta$ is chosen small enough.
\end{enumerate}

\end{document}